% Test de charge des thématiques à organismes proposés : utilisateurs simulés par HTTP.
% Compilation : latexmk -pdf test-de-charge-organismes.tex
\documentclass[11pt,a4paper]{article}
\usepackage[french]{babel}
\usepackage[T1]{fontenc}
\usepackage{lmodern,microtype,booktabs,tabularx,xcolor,hyperref}
\usepackage[margin=2.2cm]{geometry}
\newcommand{\code}[1]{\texttt{\detokenize{#1}}}
\title{Test de charge des thématiques à organismes proposés de \code{mod_stage}}
\author{Sébastien Lefebvre}
\date{17 novembre 2026}

\begin{document}
\maketitle

\section{Résumé}

Quatre vagues d'utilisateurs simulés ont été jouées contre un Moodle complet : trois promotions de
400~étudiants, avec 16, 32 puis 64~utilisateurs simultanés, et une promotion de 1\,000~étudiants
avec 48~utilisateurs simultanés. Chaque vague enchaîne l'import Excel des organismes, l'ouverture
des vœux, leur clôture,
l'appariement et le dépôt des demandes de convention, en provoquant les situations de concurrence
qu'une promotion réelle peut produire.

\begin{itemize}
  \item Aucune des 15\,959~réponses mesurées n'est une erreur. Le journal du serveur ne contient ni
    avertissement PHP ni réponse~500.
  \item Les neuf invariants relus en base après chaque vague sont respectés : aucune place
    dépassée, aucune affectation en double, aucun stage en double, aucun organisme en double.
  \item Le débit plafonne entre 30 et 33~requêtes par seconde sur la machine de test, quelle que
    soit la concurrence : le serveur est saturé, les durées s'allongent, rien ne casse.
  \item Le test a fait apparaître un défaut de performance, corrigé : la page \og Vœux et
    appariement\fg{} pesait 11~Mo pour 1\,000~étudiants (section~\ref{sec:defaut}).
  \item Une observation d'usage, sans défaut de code : quand les vœux se concentrent sur les mêmes
    organismes, des étudiants restent sans affectation alors que des places sont libres
    (section~\ref{sec:voeux}).
\end{itemize}

\section{Méthode}

\subsection{Environnement}

Moodle~4.5.16, PHP~8.3.6, PostgreSQL~16, sur une machine de 4~cœurs et 15~Go de mémoire qui héberge
aussi la base et les clients simulés. Le serveur web est celui intégré à PHP, avec 16~processus.
Les courriels sont désactivés (\code{$CFG->noemailever}) : le courriel de résultat est composé mais
pas remis. Le module est celui de la branche, version \code{2026111700}.

Ce n'est pas une configuration de production : pas de cache d'opcodes partagé entre processus,
clients et serveur sur les mêmes cœurs. Les durées se comparent entre elles, pas à celles d'un
serveur réel, qui seront plus courtes.

\subsection{Données}

\code{setup.php} crée une promotion, une thématique à organismes proposés (cinq vœux à classer,
demandes hors liste interdites), des organismes dont les places couvrent 85~\% de l'effectif, dix
enseignants référents et quatre comptes DEVE. Chaque vague a son propre site.

\subsection{Utilisateurs simulés}

\code{load.py} ouvre une session Moodle par étudiant et envoie les mêmes requêtes qu'un navigateur :
pages, formulaires avec leur clé de session, redirections. Les vœux suivent une loi de Zipf : le
premier organisme est demandé deux fois plus que le deuxième, trois fois plus que le troisième, et
ainsi de suite. C'est le cas défavorable, où presque toute la promotion vise les mêmes places.

\begin{enumerate}
  \setcounter{enumi}{-1}
  \item \textbf{Import Excel.} Un compte DEVE télécharge le modèle et le réimporte tel quel, puis
    importe un classeur qui ajoute 20~organismes et en met 10 à jour. Deux comptes DEVE importent
    au même instant le même classeur de 15~organismes nouveaux. Un classeur de 2\,000~lignes est
    importé puis réimporté dans une autre thématique. Un fichier qui n'est pas un classeur et un
    classeur sans colonne du nom sont envoyés. Les organismes ajoutés entrent dans la liste
    proposée aux étudiants.
  \item \textbf{Vœux.} Chaque étudiant ouvre \og Mes stages\fg{}, la page des organismes, puis
    enregistre ses vœux. Un sur trois les modifie ensuite, un sur dix envoie deux fois le
    formulaire en même temps, un sur vingt les retire puis les refait, un sur vingt tente trois
    envois invalides (doublon, rang vide, organisme d'une autre thématique). Pendant ce temps, un
    compte DEVE réimporte en boucle le classeur des organismes et un classeur de 40~lignes toutes
    en erreur (places illisibles, SIRET invalide).
  \item \textbf{Clôture.} Trois comptes DEVE lancent l'appariement au même instant, pendant que
    des étudiants renvoient encore leurs vœux.
  \item \textbf{Dernière place.} Dix fois de suite, une place est ajoutée à un organisme complet et
    deux comptes DEVE y affectent au même instant deux étudiants différents.
  \item \textbf{Demandes de convention.} Chaque étudiant affecté ouvre sa demande préremplie et
    l'envoie, deux fois en même temps pour un sur six. Les autres tentent une demande hors liste,
    d'abord interdite, puis autorisée : sans SIRET, puis avec SIRET ou à l'étranger.
\end{enumerate}

À partir de la phase~1, un quatrième compte DEVE recharge en boucle la liste des organismes, la
page des vœux et l'export des affectations.

Un refus prévu par l'application (vœux clos, organisme complet, SIRET exigé) est compté à part. Une
réponse est une erreur si son code est 500 ou plus, si elle contient une trace d'erreur de PHP ou
de Moodle, ou si elle n'est ni l'enregistrement attendu ni un refus prévu.

\subsection{Invariants}

Après chaque vague, \code{check.php} relit la base :

\begin{enumerate}
  \item aucun organisme n'a plus d'étudiants affectés que de places ;
  \item un étudiant a au plus une affectation par thématique ;
  \item une affectation issue de l'appariement porte sur un vœu de l'étudiant, au rang enregistré ;
  \item les vœux d'un étudiant sont au nombre demandé, de rangs 1 à 5, sans doublon ;
  \item le stage rattaché à une affectation appartient à cet étudiant et à cette thématique ;
  \item aucun étudiant n'a deux stages actifs sur la thématique ;
  \item un stage hors liste en France porte un SIRET valide ;
  \item tout stage déposé a sa demande de convention et son détail ;
  \item aucun organisme n'est proposé deux fois (même SIRET, ou à défaut même nom) dans une
    thématique.
\end{enumerate}

\section{Résultats}

\subsection{Vue d'ensemble}

\begin{center}
\begin{tabularx}{\linewidth}{@{}X r r r r@{}}
\toprule
& \textbf{Vague 1} & \textbf{Vague 2} & \textbf{Vague 3} & \textbf{Vague 4} \\
\midrule
Étudiants simulés & 400 & 400 & 400 & 1\,000 \\
Utilisateurs simultanés & 16 & 32 & 64 & 48 \\
Organismes proposés, après import & 115 & 115 & 115 & 185 \\
Places & 391 & 391 & 391 & 889 \\
Requêtes HTTP mesurées & 2\,990 & 2\,863 & 2\,967 & 7\,139 \\
Durée des cinq phases (s) & 98{,}4 & 91{,}6 & 92{,}7 & 218{,}8 \\
Débit moyen (requêtes/s) & 30{,}4 & 31{,}3 & 32{,}0 & 32{,}6 \\
Réponses en erreur & 0 & 0 & 0 & 0 \\
Invariants violés & 0 & 0 & 0 & 0 \\
Réponse la plus lente (s) & 3{,}87 & 3{,}62 & 5{,}21 & 4{,}21 \\
\bottomrule
\end{tabularx}
\end{center}

Les connexions des étudiants, faites avant chaque vague, ne sont pas mesurées.

\subsection{Durées par action}

Médiane et 95\textsuperscript{e}~centile, en secondes.

\begin{center}
\small
\begin{tabularx}{\linewidth}{@{}X r r r r r r r r@{}}
\toprule
& \multicolumn{2}{c}{\textbf{16}} & \multicolumn{2}{c}{\textbf{32}} & \multicolumn{2}{c}{\textbf{64}} & \multicolumn{2}{c}{\textbf{48, 1\,000 ét.}} \\
\textbf{Action} & méd. & p95 & méd. & p95 & méd. & p95 & méd. & p95 \\
\midrule
Page \og Mes stages\fg{} & 0{,}36 & 0{,}75 & 0{,}65 & 1{,}08 & 1{,}26 & 1{,}63 & 0{,}96 & 1{,}40 \\
Page des organismes et des vœux & 0{,}32 & 0{,}65 & 0{,}62 & 1{,}02 & 1{,}28 & 1{,}71 & 0{,}96 & 1{,}37 \\
Enregistrement des vœux & 0{,}65 & 1{,}29 & 1{,}18 & 1{,}99 & 2{,}26 & 2{,}98 & 1{,}79 & 2{,}60 \\
Vœux, double clic & 0{,}94 & 1{,}57 & 1{,}31 & 1{,}99 & 2{,}22 & 2{,}92 & 1{,}90 & 2{,}74 \\
Vœux pendant la clôture & 0{,}48 & 0{,}98 & 0{,}99 & 1{,}56 & 2{,}01 & 2{,}58 & 1{,}42 & 2{,}08 \\
Appariement (médiane de trois) & 0{,}79 & -- & 1{,}83 & -- & 2{,}56 & -- & 2{,}25 & -- \\
Affectation par la DEVE & 0{,}21 & 0{,}38 & 0{,}21 & 0{,}42 & 0{,}19 & 0{,}34 & 0{,}28 & 0{,}48 \\
Demande préremplie, page & 0{,}39 & 0{,}70 & 0{,}85 & 1{,}23 & 1{,}58 & 2{,}13 & 1{,}24 & 1{,}71 \\
Demande préremplie, envoi & 0{,}70 & 1{,}20 & 1{,}49 & 2{,}26 & 2{,}86 & 3{,}90 & 2{,}36 & 3{,}11 \\
Demande hors liste avec SIRET & 0{,}53 & 0{,}90 & 1{,}07 & 2{,}04 & 2{,}28 & 3{,}31 & 2{,}07 & 2{,}77 \\
Liste des organismes (DEVE) & 0{,}31 & 0{,}65 & 0{,}51 & 1{,}01 & 0{,}61 & 1{,}74 & 0{,}92 & 1{,}69 \\
Vœux et appariement (DEVE) & 0{,}46 & 0{,}81 & 0{,}74 & 1{,}37 & 0{,}97 & 1{,}84 & 1{,}35 & 1{,}85 \\
Export des affectations (DEVE) & 0{,}17 & 0{,}61 & 0{,}36 & 0{,}77 & 0{,}83 & 1{,}66 & 0{,}87 & 1{,}33 \\
Réimport Excel pendant les vœux (DEVE) & 0{,}77 & 1{,}06 & 0{,}99 & 1{,}21 & 1{,}70 & 2{,}13 & 1{,}46 & 1{,}92 \\
\bottomrule
\end{tabularx}
\end{center}

Le débit reste entre 30 et 33~requêtes par seconde de 16 à 64~utilisateurs simultanés : les quatre
cœurs sont occupés dès la première vague. Au-delà, les requêtes attendent leur tour et les durées croissent
avec la concurrence, sans erreur ni délai dépassé. L'enregistrement des vœux et le dépôt d'une
demande sont les actions les plus longues : elles prennent le verrou de l'activité, qui met en
file les écritures d'une même promotion. C'est ce verrou qui garantit les invariants.

L'appariement lui-même dure environ trois secondes pour 1\,000~étudiants, sous la charge des vœux
envoyés au même moment.

\subsection{Concurrence}

\begin{center}
\begin{tabularx}{\linewidth}{@{}X r r r r@{}}
\toprule
& \textbf{Vague 1} & \textbf{Vague 2} & \textbf{Vague 3} & \textbf{Vague 4} \\
\midrule
Vœux invalides refusés & 78 & 69 & 87 & 168 \\
Doubles clics sur les vœux (envois) & 114 & 82 & 98 & 222 \\
Appariements lancés au même instant & 3 & 3 & 3 & 3 \\
\quad dont ayant affecté des étudiants & 1 & 1 & 1 & 1 \\
Étudiants affectés par l'appariement & 320 & 328 & 313 & 707 \\
Vœux envoyés pendant la clôture : acceptés & 16 & 50 & 94 & 71 \\
\quad refusés, vœux clos & 32 & 46 & 98 & 73 \\
Dernière place disputée par deux DEVE & 10 & 10 & 10 & 10 \\
\quad obtenue par un seul étudiant & 10 & 10 & 10 & 10 \\
\quad refusée à l'autre, organisme complet & 10 & 10 & 10 & 10 \\
Demandes préremplies, double clic (paires) & 53 & 56 & 49 & 110 \\
Stages rattachés à une affectation & 330 & 338 & 323 & 717 \\
Affectations & 330 & 338 & 323 & 717 \\
Hors liste refusé, option décochée & 70 & 62 & 77 & 283 \\
Hors liste refusé sans SIRET & 70 & 62 & 77 & 283 \\
Hors liste accepté avec SIRET & 48 & 48 & 59 & 209 \\
Hors liste accepté à l'étranger & 22 & 14 & 18 & 74 \\
\bottomrule
\end{tabularx}
\end{center}

\begin{itemize}
  \item \textbf{Trois appariements simultanés.} Un seul affecte les étudiants ; les deux autres,
    qui passent après lui sous le verrou, ne trouvent plus personne à affecter et ne changent
    rien.
  \item \textbf{Vœux pendant la clôture.} Un envoi arrivé avant l'appariement est enregistré et
    pris en compte ; un envoi arrivé après est refusé avec le message \og vœux clos\fg{}. Aucun
    n'est perdu ni enregistré à moitié.
  \item \textbf{Dernière place.} Sur les dix courses de chaque vague, un étudiant obtient la place
    et l'autre reçoit \og organisme complet\fg{}. Jamais les deux.
  \item \textbf{Double clic sur la demande.} Chaque paire d'envois simultanés ne crée qu'un stage :
    le nombre de stages rattachés à une affectation est égal au nombre d'affectations.
  \item \textbf{Hors liste.} Toutes les demandes sont refusées tant que l'option est décochée.
    Une fois cochée, toutes sont refusées sans SIRET en France, et acceptées avec un SIRET ou pour
    un stage à l'étranger.
\end{itemize}

\subsection{Import Excel}

\begin{center}
\begin{tabularx}{\linewidth}{@{}X r r r r@{}}
\toprule
& \textbf{Vague 1} & \textbf{Vague 2} & \textbf{Vague 3} & \textbf{Vague 4} \\
\midrule
Modèle réimporté tel quel : ajoutés / mis à jour & 0/80 & 0/80 & 0/80 & 0/150 \\
Classeur mixte : ajoutés / mis à jour & 20/10 & 20/10 & 20/10 & 20/10 \\
Deux imports simultanés de 15 organismes : créés & 15 & 15 & 15 & 15 \\
Organismes en double après la vague & 0 & 0 & 0 & 0 \\
Réimports pendant les vœux & 38 & 25 & 15 & 41 \\
\quad dont ayant ajouté un organisme & 0 & 0 & 0 & 0 \\
Classeurs de 40 lignes en erreur importés & 38 & 25 & 15 & 41 \\
\quad lignes refusées et signalées & 1\,520 & 1\,000 & 600 & 1\,640 \\
Fichier illisible, colonne du nom absente : refusés & 2 & 2 & 2 & 2 \\
Import de 2\,000 lignes (s) & 2{,}15 & 1{,}78 & 1{,}71 & 1{,}69 \\
Réimport des 2\,000 lignes (s) & 3{,}87 & 3{,}34 & 3{,}34 & 3{,}10 \\
Réimport pendant les vœux, médiane / p95 (s) & 0{,}77 / 1{,}06 & 0{,}99 / 1{,}21 & 1{,}70 / 2{,}13 & 1{,}46 / 1{,}92 \\
\bottomrule
\end{tabularx}
\end{center}

\begin{itemize}
  \item \textbf{Modèle.} Le classeur écrit par Moodle se relit sans perte : réimporté tel quel,
    il met à jour chaque organisme et n'en ajoute aucun.
  \item \textbf{Imports simultanés.} Les deux imports du même classeur créent 15~organismes à eux
    deux, jamais 30 : le second, qui passe après le premier sous le verrou de l'activité, les
    reconnaît et les met à jour.
  \item \textbf{Réimports pendant les vœux.} Aucun n'ajoute d'organisme ni ne fait échouer un
    enregistrement de vœux. Leur durée suit celle des autres écritures, qui partagent le même
    verrou.
  \item \textbf{Lignes en erreur.} Chaque ligne refusée est signalée avec son numéro ; aucune
    n'est enregistrée, et le classeur ne provoque pas d'erreur.
  \item \textbf{Volume.} 2\,000~organismes s'importent en deux secondes environ. Le réimport,
    qui relit chaque organisme avant de le mettre à jour, en prend moins de quatre.
\end{itemize}

\subsection{MariaDB}\label{sec:mariadb}

Les quatre vagues ci-dessus tournent sur PostgreSQL. Une cinquième a été jouée sur MariaDB~10.11,
la version de la CI, avec un second Moodle : 400~étudiants, 32~utilisateurs simultanés, les mêmes cinq
phases. MariaDB isole les transactions autrement que PostgreSQL et Moodle y pose ses verrous par un
autre mécanisme : c'est là qu'un défaut de concurrence propre à cette base se verrait.

\begin{center}
\begin{tabularx}{0.8\linewidth}{@{}X r@{}}
\toprule
\textbf{Vague MariaDB} & \\
\midrule
Requêtes HTTP mesurées & 2\,870 \\
Réponses en erreur & 0 \\
Invariants violés & 0 \\
Organismes en double après les imports & 0 \\
Appariements simultanés ayant affecté des étudiants & 1 sur 3 \\
Étudiants affectés par l'appariement & 321 \\
Dernière place : obtenue par un seul étudiant & 10 sur 10 \\
Stages rattachés à une affectation / affectations & 331 / 331 \\
Débit moyen (requêtes/s) & 54{,}7 \\
Réponse la plus lente (s) & 1{,}51 \\
\bottomrule
\end{tabularx}
\end{center}

Le comportement est le même que sur PostgreSQL : un seul appariement, une seule affectation par
place disputée, aucun stage ni organisme en double. Le débit, plus élevé, ne se compare pas à celui
des autres vagues : ce second site, tout juste installé, est plus léger que le premier.

Le journal du serveur contient une réponse~500, sur le tableau de bord de Moodle
(\code{/my/}), pendant les connexions simultanées qui précèdent la vague. Elle ne vient pas du
module et n'entre pas dans les mesures.

\section{Défaut corrigé}\label{sec:defaut}

La page \og Vœux et appariement\fg{} proposait, sur la ligne de chaque étudiant, un menu de tous
les organismes. Son poids croissait comme le produit des deux nombres.

\begin{center}
\begin{tabularx}{0.8\linewidth}{@{}X r r@{}}
\toprule
\textbf{Promotion} & \textbf{Avant} & \textbf{Après} \\
\midrule
400 étudiants, 80 organismes & 2{,}63~Mo & 0{,}58~Mo \\
1\,000 étudiants, 150 organismes & 11{,}32~Mo & 1{,}23~Mo \\
\bottomrule
\end{tabularx}
\end{center}

Le menu de chaque ligne se limite désormais aux vœux de l'étudiant et à son affectation actuelle.
Un formulaire unique, \og Affecter hors vœux\fg{}, en haut de page, permet d'affecter n'importe
quel étudiant à n'importe quel organisme. Les quatre vagues ci-dessus ont été jouées après cette
correction, puis rejouées après les corrections de l'audit du code
(\code{audit-organismes-proposes.tex}) ; leurs résultats sont ceux du code livré.

\section{Observation sur les vœux}\label{sec:voeux}

\begin{center}
\begin{tabularx}{\linewidth}{@{}X r r r r@{}}
\toprule
& \textbf{Vague 1} & \textbf{Vague 2} & \textbf{Vague 3} & \textbf{Vague 4} \\
\midrule
Étudiants affectés par l'appariement & 320 (80~\%) & 328 (82~\%) & 313 (78~\%) & 707 (71~\%) \\
\quad sur leur premier vœu & 114 & 118 & 132 & 231 \\
\quad sur leur deuxième ou troisième vœu & 132 & 145 & 122 & 314 \\
\quad sur leur quatrième ou cinquième vœu & 78 & 68 & 61 & 163 \\
Étudiants sans affectation & 80 & 72 & 87 & 293 \\
Places restées libres après l'appariement & 71 & 63 & 78 & 182 \\
\bottomrule
\end{tabularx}
\end{center}

Les rangs sont relevés en fin de vague : ils comptent aussi les quelques affectations faites à la
main en phase~3 lorsqu'elles tombaient sur un vœu de l'étudiant.

Avec des vœux très concentrés, un cinquième à près d'un tiers des étudiants restent sans affectation alors que
des places sont libres dans des organismes que personne, ou presque, n'a classés. Ce n'est pas un
défaut de l'appariement : il ne peut affecter un étudiant qu'à un organisme qu'il a demandé. La
concentration simulée est volontairement forte ; une promotion réelle sera sans doute moins
tranchée, mais l'effet existe dès que quelques organismes attirent la majorité des vœux.

Trois leviers, tous disponibles : augmenter le nombre de vœux à classer, rouvrir les vœux pour un
tour complémentaire sur les places restantes, ou affecter à la main les derniers étudiants.

\section{Limites}

\begin{itemize}
  \item La machine de test porte à la fois le serveur, la base et les clients. Les durées
    absolues ne valent pas pour un serveur de production.
  \item MariaDB n'a reçu qu'une vague, de 400~étudiants (section~\ref{sec:mariadb}) ; les quatre
    autres, dont celle de 1\,000~étudiants, tournent sur PostgreSQL.
  \item Les clients n'exécutent pas le JavaScript des pages et ne chargent pas les feuilles de
    style : seul le travail du serveur est mesuré.
  \item La reprise de stages précédents, action ponctuelle d'un seul compte DEVE, ne fait pas
    partie de la vague. Elle est couverte par les tests unitaires.
  \item Les classeurs importés sont écrits par \code{openpyxl} et par Moodle, pas par Excel ni
    LibreOffice.
  \item L'envoi réel des courriels de résultat n'est pas mesuré.
\end{itemize}

\section{Tests unitaires}

La fonction est couverte par 46~tests PHPUnit (\code{tests/placement_test.php}), auxquels
s'ajoute le test de sauvegarde et de restauration, et par quatre scénarios Behat
(\code{tests/behat/placement.feature}). Les deux suites PHPUnit passent sur PostgreSQL~16 et sur
MariaDB~10.11. Ils vérifient notamment les états de la
campagne, la validation des vœux, l'ordre et les places dans l'appariement, l'équité du tirage au
sort sur 400~tirages, une promotion entière de 300~étudiants, le tour complémentaire, la
réouverture, l'affectation par la DEVE, la règle du hors liste et du SIRET dans les deux
formulaires de demande, le préremplissage, l'import Excel, la reprise de stages précédents, la
tâche planifiée et son courriel, l'export, l'encart de l'étudiant, la réinitialisation, l'import
de thématiques et les données personnelles.

\section{Reproduire}

Les scripts et leurs journaux sont dans \code{docs/audit/organismes/} ; le fichier
\code{README.md} du dossier donne les commandes.

\end{document}
